% ==========================================================================
\chapter{核心代码与复现说明}
% ==========================================================================

本附录给出论文核心代码的结构与关键片段。完整代码位于配套的 \texttt{Code} 文件夹（其下的 \texttt{Undergraduate Thesis Code} 子文件夹），数据位于 \texttt{data} 子文件夹，图片位于 \texttt{figures} 子文件夹。

\section{代码结构}

\begin{table}[htbp]
  \centering
  \small
  \caption{代码文件与功能}
  \label{tab:code-files}
  \begin{tabularx}{\textwidth}{lX}
    \toprule
    文件 & 功能 \\
    \midrule
    \texttt{burgers\_imex.py} & 有限体积离散、Rusanov 通量、MUSCL/minmod、IMEX-Euler \\
    \texttt{stability\_analysis.py} & 双变量稳定性函数、圆盘/相图、联合稳定域、离散符号 \\
    \texttt{run\_burgers\_experiments.py} & 刚性相图、最大步长、激波精度、收敛性 \\
    \texttt{plot\_results.py} & 从缓存数据绘制论文图 \\
    \texttt{jinxin.py} & Jin--Xin 松弛系统、两种源项分裂、谱分析 \\
    \texttt{run\_jinxin\_experiments.py} & 谱半径扫描、临界 CFL、数值对比 \\
    \texttt{stiff\_ode\_example.py} & 第二章刚性 ODE 算例 \\
    \bottomrule
  \end{tabularx}
\end{table}

\clearpage
\section{IMEX-Euler 的时间推进}

下面的片段对应 \texttt{burgers\_imex.py} 中的核心循环。矩阵 $\bA=\bI-(1-\theta)\nu\Delta t\bD$ 只分解一次，随后每步用三角回代求解。

\begin{lstlisting}[language=Python,caption={IMEX-Euler 时间推进的核心循环}]
c = (1.0 - theta) * nu * dt / dx**2
main = np.full(N, 1.0 + 2.0 * c)
off = np.full(N - 1, -c)
A = diags([off, main, off], [-1, 0, 1], format="csc")
lu = splu(A)

for n in range(nsteps):
    conv = rhs_convective(u, grid, bcl, bcr, order=order)
    lap = laplacian(u, grid, bcl, bcr)
    rhs = u + dt * (conv + theta * nu * lap)
    rhs[0] += c * bcl
    rhs[-1] += c * bcr
    u = lu.solve(rhs)
\end{lstlisting}

\section{双变量放大因子}

\begin{lstlisting}[language=Python,caption={离散符号与放大因子}]
def amplification(xi, dt, nu, theta, a, dx):
    lam_c = (-1j * a * np.sin(xi) / dx
             - abs(a) * (1.0 - np.cos(xi)) / dx)
    lam_d = -2.0 * (1.0 - np.cos(xi)) / dx**2
    z_e = dt * (lam_c + theta * nu * lam_d)
    z_i = dt * (1.0 - theta) * nu * lam_d
    return (1.0 + z_e) / (1.0 - z_i)
\end{lstlisting}

\section{Jin--Xin 源项更新}

\begin{lstlisting}[language=Python,caption={策略 A 与策略 B 的源项更新}]
du, dv = transport_rhs(u, v, grid.dx, a, uL, uR)
u_new = u + dt * du
if strategy == "A":
    v_new = (v + dt * dv + mu * f_burgers(u_new)) / (1.0 + mu)
elif strategy == "B":
    v_new = (v + dt * dv + mu * f_burgers(u)) / (1.0 + mu)
\end{lstlisting}

\section{谱半径扫描}

\begin{lstlisting}[language=Python,caption={策略 B 刚性极限的谱半径}]
g = cfl * (1.0 - np.cos(xi))
h = (cfl / a) * np.sin(xi)
M = np.array([[1.0 - g, -1j * h], [b, 0.0]])
rho = np.max(np.abs(np.linalg.eigvals(M)))
\end{lstlisting}

\section{复现步骤}

代码依赖 Python 3.10 以上版本、NumPy、SciPy 与 Matplotlib。为避免污染系统环境，可将依赖安装到本地目录：

\begin{lstlisting}[language=bash,caption={安装依赖}]
cd "Code/Undergraduate Thesis Code"
python3 -m pip install --target .pylibs \
    --break-system-packages numpy scipy matplotlib
\end{lstlisting}

然后依次运行：

\begin{lstlisting}[language=bash,caption={生成图片与数据}]
cd "Code/Undergraduate Thesis Code"
export PYTHONPATH=.:.pylibs
export MPLCONFIGDIR=/tmp/mplconfig
python3 stability_analysis.py
python3 run_burgers_experiments.py all
python3 run_jinxin_experiments.py
python3 plot_results.py
\end{lstlisting}

全部实验在普通笔记本电脑上约需数分钟。随机数未使用，结果可完全复现。
